\chapter{How grammar feels} \label{ch:How grammar feels}

%\begin{quote}
% On the modest explanation, linguistic intuitions are, like any other intuitive judgments a person may make, immediate and unreflective reactions formed in a central processing system.
%\end{quote} \citep[4 (introduction)]{Schindler2020a}

%\begin{quote}
% ``I may long for certainty, but I have to live with doing the best I can. I may concoct a myth to explain that my certainty, unlike yours, taps into universal moral truths. Reality will soon dissolve that myth. Voltaire (1694–1778), a French philosopher of the Enlightenment, concisely summed up the state of affairs: \textit{Uncertainty is an uncomfortable position, but certainty is an absurd one.''\textsuperscript{10} Voltaire is right, of course, but nevertheless, we must make some choice or other, and not acting can be the worst of the options. Knowing that I am doing the best I can seems like cold comfort, yet my inner Socratic voice recommends that I make do with cold comfort.''\cite[Introduction]{churchland2019}
%\end{quote}

When you encounter an ungrammatical utterance, what does it feel like? Is it akin to stepping off the sidewalk unexpectedly~-- a sudden jolt to your cognitive equilibrium? Or sometimes more than just a simple misstep~-- a momentary disorientation, where the flow of language and thought breaks?

Is it sand in your shoe, a branch that scratches you as you walk by, an ice-cream headache, the sting of a mosquito, the glare of headlights? Does it feel heavy or light, salty or sour? Does it have a colour?

%Does ungrammaticality lean towards the cool side, like a draft that sneaks in through an unseen gap. Or is it a mouthful of scalding coffee? Is it a path that leads to a dead end or a door that's slightly ajar but won't fully open?

Emotionally, is this encounter more likely to evoke irritation and discomfort, fear and anger, or dismissal and superiority?

Or can it be all of those things at different times?

\bigskip

The idea that we have only five senses~-- sight, smell, hearing, taste, and touch~-- is nonsense. The traditional list of five ignores pain, pressure, itch, temperature, and chemical irritants like acids, the capsicum in hot peppers, and allyl isothiocyanate in radishes, wasabi and mustard. And that's just the exteroceptive senses~-- those that detect stimuli from outside the body.

We also have a rich array of interoceptive senses monitoring the internal states of our bodies, hunger for example, or sleepiness. When someone is dizzy, we literally say they've lost their sense of balance. Our sense of self relies on proprioception, the feeling we have of our bodies, their coherence, their locations, and their movements in space. We have senses of thirst, nausea, and fullness. We even have a sense of the levels of carbon dioxide and oxygen in our blood. We're sensitive creatures.

Sensations arise from the interaction between internal or external stimuli, our sense organs, and our nervous system. When our elbow is bent at 90$^{\circ}$, intrafusal fibers in our muscles register the relative amount of stretch in our biceps and triceps and Ruffini endings in the joint capsule and ligaments detect changes in joint angle. Nerves pick up the signals from these mechanoreceptors, and the sensations are sent as proprioceptive signals to our brain, informing it of the arm's position and movement. But positional sensations like this are typically evaluatively neutral. 90$^{\circ}$ should be neither more nor less pleasant than 120$^{\circ}$.

When we do evaluate sensations as good or bad, we've moved on to talking about \textsc{feelings}, the positive or negative experiences that come from sensing the world and ourselves. We sense hunger, pain, or high levels of CO\textsubscript{2} in the blood and we feel that they are unpleasant, problems to be righted. The sensation of warm touch on the arm, the sweetness of a peach, and the temporary relief that comes from scratching a mosquito bite, in contrast, foster positive feelings of relief, warmth, and connection, directing us to seek out such experiences.

\begin{quote}
    No sooner had the warm liquid mixed with the crumbs touched my palate than a shudder ran through me and I stopped, intent upon the extraordinary thing that was happening to me. An exquisite pleasure had invaded my senses, something isolated, detached, with no suggestion of its origin. And at once the vicissitudes of life had become indifferent to me, its disasters innocuous, its brevity illusory this new sensation having had on me the effect which love has of filling me with a precious essence; or rather this essence was not in me it was me. ... Whence did it come? What did it mean? How could I seize and apprehend it? ... And suddenly the memory revealed itself. \par\raggedleft~-- Marcel Proust, \textit{In search of lost time}
\end{quote}

Proust's vivid description highlights how a simple sensory encounter~-- the taste of a madeleine dipped in tea~-- can be intensely pleasurable. But the experience transcends his sensory delight. He perceives a feeling whose source he can't identify at first, one triggered by but distinct from his experience of the smells, flavours, and textures of the moistened cake. Proust recognizes an upwelling associated with this particular set of sensations even before the memory arrives. It announces itself the way a yawn does, with a gradual build-up of tension before the mouth opens wide and the breath begins to move. A pricking of the thumbs. Summer's ripening breath.

This kind of feeling is called a metacognitive feeling. Unlike the taste of a madeleine or the warmth of a cup of tea, metacognitive feelings are generated by our own thoughts and mental operations. Such impressions include the positive feeling of knowing you might experience when reviewing your responses on a test, the pleasant recognition of a friend's familiar gait, or the empowering conviction of the rightness of your opinion in a debate. It's not the thoughts that we have but the feeling that we have thoughts, whether we can access them or not. When Faulkner writes, \enquote{memory believes before knowing remembers,} he's evoking metacognitive feelings.

Metacognitive feelings aren't all positive. Tip-of-the-tongue feelings bedevil us. The thought that you've left something undone can be distressing. And then there's the feeling that this book is all about: ungrammaticality. This too I would place among our metacognitive feelings.


%Nick Chater, in \textit{The Mind is Flat}, argues that mental processes are fundamentally superficial, challenging the deep and complex models often posited in psychology and cognitive science. From this perspective, metacognitive feelings, including the sensation of (un)grammaticality, are immediate and surface-level responses, not rooted in any deep, structured understanding of grammar or cognitive processing.

%Chater would likely suggest that the feeling of (un)grammaticality arises not from an intricate, rule-based analysis of language structures internalized through extensive learning, but from heuristic responses generated by the brain's superficial processing systems. These systems are quick to judge based on patterns and examples encountered in everyday language use, rather than deep grammatical rules. For instance, a sentence might feel ungrammatical simply because it's unfamiliar or deviates from frequently observed patterns, not necessarily because it violates a specific grammatical rule.

%This approach emphasizes how our perceptions of grammaticality are shaped more by pragmatic cues and less by any robust internal grammar engine. Thus, the feeling of (un)grammaticality is a product of our minds navigating through language based on surface impressions and immediate judgments, aligning well with Chater's view that the mind constructs its seemingly complex operations from moment to moment, without any underlying depth.

%\bigskip

%But we \emph{do} have a sense of grammar, or at least something it feels like to encounter grammaticality and ungrammaticality~-- not one thing, just like pain isn't one thing, but a set of related feelings that we experience when we encounter grammatical and ungrammatical sentences and which we associate with grammaticality and ungrammaticality.

When we identify a sentence as ungrammatical, our response isn't primarily rooted in a physical sensation but in a cognitive evaluation. This process involves evaluating incoming linguistic constructions to ensure that they're familiar, conventional, and congruent, that they evoke the appropriate meaning for the situation, and that, if not, they deviate in understandable ways to good effect. It's difficult to identify necessary and sufficient criteria for feelings of ungrammaticality, but cognitive dissonance is often at their heart.

%This cognitive evaluation of grammaticality is reflective of our linguistic experience, a complex product of language exposure, learning, and use. It is through this lens that we can understand our reactions to ungrammatical constructions not as visceral responses but as cognitive acknowledgements of variance from the norm. The experience of encountering ungrammatical language, therefore, speaks to our capacity to detect, process, and interpret deviations within the framework of our linguistic knowledge.

%In this context, grammaticality is less about provoking a sensory reaction and more about the brain's ability to process language in accordance with learned patterns, what Dan Dennett called real patterns. %How do I connect this to chapter 10? Shall I just deal with it here?
Our reactions to grammatical incongruences may vary widely, from irritation to curiosity, reflecting the diversity in our linguistic backgrounds and tolerances for variation. Ultimately, the sense of grammar underscores the cognitive foundations of language processing and our adaptive strategies for navigating linguistic complexity.

%Is your sense of ungrammaticality the same feeling of wrongness you get when you see a stone levitating? I think it evokes more curiosity too, though. What about the sense of wrongness you feel when seeing unfairness?

%How do you react to your own grammatical slips?

%Perhaps it would be a step to far, though, to say we have a sense of grammar. I'll settle for saying that there is something it feels like to encounter an ungrammaticality~-- not one thing, just like pain isn't one thing, but a set of related feelings that we experience when we encounter ungrammatical sentences and which we associate with ungrammaticality.

\section{Whose gorilla?}

In 1973, Jorge Hankamer and Paul Postal published a short squib called \enquote{Whose gorilla?} \citep{Hankamer1973}. Its joke was that English seemed to have a hole in its pronoun system. You could ask \mention{whose gorilla is this}, and you could ask \mention{whose is this}, but you seemingly couldn't use relative \mention{whose} on its own: *\mention{the guy whose you saw}.

The full story of this missing-but-not-missing \mention{whose} now belongs to Chapter \ref{ch:getting-grammaticality-wrong}. Here, it matters because the case shows what a grammatical feeling is: immediate, compelling, and sometimes wrong.


\section{Free adjuncts, danglers, and howlers}

Dangling participles haunt the dreams of grammarians. \textit{\uline{Turning his attention back to his work}, a dangler stared him in the face.} As you first read this, you develop the idea that somebody was turning his attention back to his work, and, right after the comma you see \textit{a dangler}, and in your developing sense of the meaning, you entertain the idea that that dangler was that person. And then you encounter \textit{him}, which forces a complete reassessment of the situation. This is the experience of reading a bad dangling modifier, or what linguists call \textsc{free adjuncts}.

Consider (\ref{ex:erection}), noticed by my friend, Jim \citet{Donaldson2020}.

\ea[*]{\textit{Suddenly, he sits up and tugs my panties off and throws them on the floor. Pulling off his boxer briefs, his erection springs free. Holy cow!} \\ \hfill(E. L. James \textit{50 Shades of Grey})}\label{ex:erection}
\z

This dangler is bad. It qualifies as a howler, and, in response, I've marked it with an asterisk. I think it's ungrammatical, but is it?
% EDITORIAL-HPC: This is a prime place to separate grammaticality proper from
% discourse/processing/style. A howler may trigger the detector without failing
% the same kind of form-value coupling as agreement or complement selection.

% TODO: source-grounding -- bibliographic data in body text; move to \citep{Lovinger2000} and add bib entry
The first example cited by The Penguin Dictionary of American English Usage and Style (by Paul W. Lovinger; New York: Penguin Reference, 2000) is as clear a case as one could want, taken from a book about cannabis (if you could steel yourself, please, I want no politically incorrect giggling at this one):

\ea
 \ea[]{\textit{Although widely used by the men, Bashilange women were rarely allowed to smoke cannabis.}}
 \ex[]{\textit{Rich and creamy, your guests will never guess that this pie is light.}}
 \z
\z

\begin{quote}
\enquote{Does this fall under the no dangling modifier prescription?}, asks Rosanne, in a post on The X-bar. Yes, Rosanne, it does. Like participles, adjectives and also some idomatic preposition phrases, when used as adjuncts, need an understood subject (or, it might be better to say, a target of predication) to be filled in if they are to be understood. The prescriptive tradition says that the subject filled in must be the one obtained from the subject of the matrix clause. Here that would be your guests, which makes a nonsense reading, so the sentence cited would be treated as an error.

But the prescriptivists have a problem. Sentences of this kind, which call for you to fill in the understood subject from somewhere else (here, the subject of is light in the subordinate clause), are so common that when I and several friends have spent some time picking new ones up from print and radio sources, we get them at a rate of as many as one per day. That's in edited sources, where grammatical errors have almost entirely been screened out. This just can't be syntactic error. It's too frequent.\hfill\citep{LanguageLog218}
\end{quote}

% TODO: source-grounding -- "delves into the intricacies / nuanced perspective / challenges traditional views" is heavy AI vocabulary; cluster-cite of four LanguageLog posts is unusually generic; verify the four posts say what's attributed
Donaldson's PhD dissertation, supervised by Geoff Pullum, examines dangling modifiers in detail. \citet{LanguageLog218,LanguageLog1938,LanguageLog1973,LanguageLog2153} and Donaldson argue that while danglers can be seen as \enquote{ill-written and discourteous}, \enquote{inconsiderate}, \enquote{stunningly inept}, or \enquote{bad grammatical manners}, they don't necessarily constitute ungrammaticality.


The construction means that his erection pulled off its boxer briefs, but that's obviously not the intended meaning. \enquote{Just about every rhetoric, grammar, and handbook written since the latter part of the 19th century warns the student against such constructions} \citep[314]{Gilman1989}.

Consider the sentence: \textit{Having arrived late, the meeting was already over.} Traditional grammar would flag this as a dangling modifier since the subject of the participle phrase \textit{Having arrived late} doesn't logically correspond to the subject of the main clause, \textit{the meeting}. Pullum and Donaldson's perspective invites us to reconsider whether this incongruence automatically renders the sentence ungrammatical or merely less clear or elegant.

Their argument suggests that the judgment of grammaticality involves more than a binary distinction between right and wrong. It encompasses considerations of clarity, elegance, and the speaker or writer's intent. From this viewpoint, the use of a dangling modifier might reflect stylistic choices or conversational shortcuts rather than a lack of grammatical competence.

Reflective equilibrium offers a valuable framework for navigating this debate. It encourages us to balance our intuitive disapproval of danglers with the understanding that language flexibility often accommodates such constructions in everyday usage. This balancing act involves weighing the clarity and communicative effectiveness of language against the descriptive realities of how people speak and write.

Examining the acceptance of dangling modifiers in different contexts and registers can shed light on the fluid boundaries of grammaticality. For instance, creative writing and informal speech frequently employ danglers to convey a sense of immediacy or emotional resonance, prioritizing expressive impact over strict adherence to grammatical rules.

\section{What does (un)grammaticality feel like?}

Sensations arise from the interaction between stimuli, our sense organs, and our nervous system. We've got the traditional five senses: sight, taste, smell, hearing, and touch. But it goes well beyond that. We have a sense of balance, regulated by nerves in our vestibular system in our inner ears. Our sense of hunger is triggered by nerves in the gastrointestinal tract and transmitted through the vagus nerve to the hypothalamus. We sense the location of our various body parts and the angles of our joints, pain, pressure, itch, hot and cold, and chemical irritants like capsicum and allyl isothiocyanate in wasabi and mustard. We even have a sense of the levels of carbon dioxide and oxygen in our blood.

We're sensitive creatures, but we have no sense organs for grammar, so grammar isn't something we sense, and (un)grammaticality isn't a sensation. But it's a feeling, or at least part of it is. Feelings are sensations that have an evaluative aspect.

\begin{quote}
    No sooner had the warm liquid mixed with the crumbs touched my palate than a shudder ran through me and I stopped, intent upon the extraordinary thing that was happening to me. An exquisite pleasure had invaded my senses, something isolated, detached, with no suggestion of its origin. And at once the vicissitudes of life had become indifferent to me, its disasters innocuous, its brevity illusory~-- this new sensation having had on me the effect which love has of filling me with a precious essence; or rather this essence was not in me it was me. ... Whence did it come? What did it mean? How could I seize and apprehend it? ... And suddenly the memory revealed itself. \\\hfill~-- Marcel Proust, \textit{In Search of Lost Time}
\end{quote}

Proust describes a feeling triggered by but unconnected to any immediate sensation, a feeling about a memory but without, at least initially, access to the memory itself. It's a feeling about a feeling, a metacognitive feeling.

Language understanding is another such process we can have feelings about, and like Proust's reverie, those feelings come in different shapes and intensities. Consider the sentences in (\ref{ex:green-cheese}). All ungrammatical, but each elicits a different reaction. The first might cause a double-take~-- what does it mean for cheese to be made of green ideas? The second feels more jarringly wrong, with words simply in the wrong places. The third is word salad, scrambled past the point where you can even try to make sense of it.


\ea\label{ex:green-cheese}
    \ea[*]{\textit{The cheese green ideas sleep late furiously.}}
    \ex[*]{\textit{Late the ideas sleep furiously green cheese.}}
    \ex[*]{\textit{Sleep furiously late ideas the green cheese.}}
    \z
\z

Several things shape the felt severity: how much of the structure is broken, what level the breakdown is at (a semantic mismatch feels different from a morphological one), and how surprising the error was given what came before.

But across all the variability, one thing stays constant: ungrammatical sentences evoke a feeling, and grammatical ones rarely do. A sentence that hews to syntax~-- and almost all the ones we encounter do~-- doesn't elicit a feeling of rightness to match the wrongness of its ungrammatical counterpart. Grammaticality is the baseline. It just feels normal.

Exceptions exist. A mellifluous line of poetry, a sharp witticism, a sentence whose clarity makes a hard idea easy~-- these can provoke a feeling. So can a well-placed correction, for those who care about that sort of thing. But for the most part, grammaticality is like clean air or a healthy back: you don't notice it until it's gone. It's homeostasis. Only when something perturbs it~-- a syntactic stumble, a semantic contradiction, a morphological mismatch~-- do we feel the friction.

So why this asymmetry? Because the feeling of ungrammaticality is a metacognitive feeling, and metacognitive feelings only register when the underlying process is in trouble.

Consider some other members of the family. The \emph{feeling of knowing} is what tells you, before you can say it, that you know who played Ron Weasley in the Harry Potter films. The \emph{tip-of-the-tongue} state is what tells you that you almost know it, but can't quite get to it. The \emph{feeling of confidence} is what guides your hand on a multiple-choice test. Each is a feeling about a cognitive process that you can't see directly: memory retrieval, lexical access, judgement under uncertainty. Pablo \textcite{FernandezVelasco2021} adds disorientation to the list. Disorientation is the feeling that rides on top of navigation. When navigation goes well, no feeling. When something goes wrong with it, you feel lost.

The same shape fits the feeling of ungrammaticality. The process beneath is parsing, and parsing works by prediction. As you read or hear a sentence, your brain runs ahead at every level~-- phonological, lexical, syntactic, semantic, pragmatic~-- predicting what comes next \citep{Clark2013,Kuperberg2016}. When the input matches the prediction, the process keeps moving and there's nothing for consciousness to do. When the prediction misses, the resulting error gets weighted, propagated, and registered \citep{Rabovsky2018}, and what surfaces in awareness is the feeling that something has gone wrong with the language.

\bigskip

The same asymmetry appears in music. Wagner's Tristan chord, F, B,
D\(\sharp\), and G\(\sharp\), from \enquote{Tristan und Isolde}, gave
audiences a sonority they had no settled way to hear. It was
unpredictable, and so it registered as wrongly dissonant, the musical
equivalent of ungrammatical. But some listeners were able to find a
sense in it, a narrative sense rather than a harmonic one.

Later, I'll get into how constructions and other linguistic forms can
mean things other than denotations and predications, things like
identity, status, deference, and interrogation. The point here is that
meaning isn't a single thing. Because some people could see analogies
between unfulfilled harmonic tension and the unrequited love central to
the opera's story, the chord wasn't dismissed.

Almost certainly it was the case that people had to work to find such a sense, and they did so because the author of the chord was Wagner. This is crucial context. Wagner had come to be seen as a musical genius, and his previous operas~-- notably Der fliegende Holländer (1843), Tannhäuser (1845), and Lohengrin (1850)~-- had featured increased chromaticism and a more flexible approach to tonality. He was also renowned for his concept of the Gesamtkunstwerk, or ``total work of art,'' which sought to unify all artistic elements within an opera. So, while the musical context of the Tristan chord evoked a sense of wrongness, this sense may have been attenuated by the context of having Wagner's name attached to it.

A lesser musician~-- someone young and unknown, or one with other such prejudicial strikes against them~-- may have been seen as simply incompetent had they composed or performed such a thing. But, although the chord didn't make sense, nor did it make sense to believe that the great Opernkomponist had made an error such as this. It had to be intentional, and so the task these listeners set themselves was to determine the intention. And they may have set off on that journey because the predictive brains had dampened the feeling of dissonance for these folks because, well, ``it's Wagner''?

Whether my stylized story is true or not, what's indisputable is that the opera continued to be performed. While it wasn't an immediate sensation, it gradually gained recognition and acclaim for its innovative music and dramatic storytelling.

The opening run itself wasn't extended, but that was common for opera premieres at the time. Subsequent productions in other cities helped to build its reputation. By the 1870s and 1880s, ``Tristan und Isolde'' had become a staple of the operatic repertoire and was widely regarded as a masterpiece. And, as a result, more and more people were exposed to the chord. More composers started using it, exposing even more audiences, and folks got used to it. It came to reside in the predictable, and it came to sound no wronger than a minor chord in the right contexts.

\bigskip

So the feeling of ungrammaticality is the conscious side of a parser surprised by what came in. Like the tip-of-the-tongue or disorientation, it isn't the cognitive work itself~-- it's the signal that the work is in trouble. And because we don't feel the work when it goes well, grammaticality has no felt counterpart. Smooth processing leaves no trace. We don't feel the absence of trouble. We only feel trouble.

\bigskip

This account also explains why grammaticality judgments vary so much across speakers. If the feeling tracks parsing, and parsing draws on whatever statistical regularities a speaker has been exposed to, then speakers with different linguistic backgrounds will have different intuitions about what's grammatical.

Two cases have made this hard to ignore for me. My daughter Mia can use \mention{they} for a named roommate without any hitch at all. I understand the intended meaning, and I endorse the change, but the older parse still fires first: for a moment, I hear \mention{they} as plural. For Mia, the sentence is ordinary. For me, it still catches.

The doubled \mention{is} has the same shape from the other direction. I hear polished speakers say things like \mention{the problem is, is that}, and the conversation barely slows. For me, the second \mention{is} still raises a flag, because it has no subject of its own. The construction may be grammatical for the people who use it. My feeling hasn't caught up with theirs.

For example, double negatives such as \textit{I don't know nothin about that} or verb forms such as \textit{I seen him yesterday} are familiar to most English speakers, but they're not part of Standard English. As a result, in a context where Standard English prevails, they may provoke a feeling of ungrammaticality. But when the Kid Rock laments, ``I ain't seen the sun shine in 3 damn days'' or Lil West boasts, ``I don't take no L's,'' we don't expect Standard English, and no feeling of ungrammaticality is evoked.

Young children over-regularize past tenses. So, instead of \textit{I found you,} a child might say, \textit{I finded you}. On the one hand, this is ungrammatical and we all know it. On the other, it doesn't feel as ungrammatical as if I said to you \textit{I wrote two book}. Even if we hadn't expected the child to say \textit{finded}, the reason for it is readily apparent in the context, and our sense of ungrammaticality is commensurately dampened. Similarly, should we hear someone with a noticeable foreign accent say, \textit{I must to go,} it registers as ungrammatical, but at the same time, the level of surprise is attenuated by the inference that this is someone who has come to English later in their life.

Conversely, I speak Japanese with a reasonable level of conversational fluency, but if a Japanese speaker confidently said (\ref{ex:fond-of-cats}b) to me instead of (a) to explain why she had a cat, I probably wouldn't notice or experience any feelings of ungrammaticality. In (a), \textit{neko} `cat' is the topic~-- thus \textsc{top}~-- and the lack of a subject suggests that a first-person subject `I'. \textsc{Cop} means a copula, roughly \textit{be}, so in a stilted overly direct translation it's `as for cats, I'm fond (of them).'

\ea\label{ex:fond-of-cats}
    \ea[]{\textit{猫~~~~~} \textit{は~~~} \textit{好き} \textit{です}。\\
      \gll Neko wa suki desu. \\
           cat \textsc{top} fond \textsc{cop} \\
      \glt `I'm fond of cats.'
      }
    \ex[]{\textit{猫~~~~~} \textit{が~~~} \textit{好き} \textit{です}。\\
      \gll Neko ga suki desu. \\
           cat \textsc{subj} fond \textsc{cop} \\
      \glt `Cats are fond of it.'
      }
    \z
\z

The second sentence, though, has \textit{cat} as the subject (\textsc{subj}), with no particular topic stated. That means the topic is probably whatever we've been talking about, and the sentence tells us that cats are fond of whatever it is. I know this \textit{ga}--\textit{wa} distinction, and I usually make the right choice when speaking. But English doesn't mark topics vs subjects this way, and so it might just slip by me if I wasn't primed for it.

\bigskip

The sense of ungrammaticality we feel isn't a binary, all-or-nothing affair. It's a graded experience, one that spans highs and lows, that reflect the uncertain nature of prediction. A minor syntactic slip, something very close to expectations, might generate a small, localized prediction error, experienced as a fleeting sense of discomfort. A more serious grammatical deviation might generate a larger, more widespread error signal, experienced as a visceral sense of linguistic wrongness or confusion.

And yet what will register as serious isn't easy to pin down. Subject--verb agreement errors seem objectively minor, often rectifiable through the mere toggling of a single \textit{--s}, but \uline{this error feel} egregious. Other times, such disagreements might slide by with hardly a niggle. Perhaps you sensed the infelicity in the previous paragraph: \textit{reflect} should have been \textit{reflects}. If you did, I'd guess the intensity of ungrammaticality you felt was minor compared with \textit{this error feel\dots}.

\subsection{Grammar sensitivity in animals}

The sensation of grammaticality we experience is often taken as evidence for the human-specific nature of language. But could this feeling have emerged from cognitive abilities that predate human language? Recent research with non-human primates suggests that the mental machinery underlying our sensitivity to grammatical dependencies may have evolved long before language itself.

In a groundbreaking study, \citet{Ravignani2013} demonstrated that squirrel monkeys (\textit{Saimiri sciureus})~-- New World primates whose lineage diverged from ours at least 36 million years ago~-- can detect abstract, non-adjacent dependencies in auditory sequences. The monkeys were habituated to tone sequences following an AB$^n$A pattern, where the first and last elements matched while being separated by a variable number of intervening elements. Not only did the monkeys distinguish between sequences containing dependencies and those lacking them, but they also generalized to previously unheard pitch classes and novel dependency distances.

This is remarkable precisely because squirrel monkeys lack vocal learning abilities and possess a relatively fixed vocal repertoire. These monkeys will never speak a human language, yet they possess the cognitive toolkit to recognize abstract patterns that form the foundation of many linguistic structures. As the authors note, ``Although no squirrel monkey will probably ever speak a human language, these monkeys possess the cognitive potential to recognize the rule generating plurals of Turkish nouns, or many other linguistic phenomena'' \citep[4]{Ravignani2013}.

What makes this finding so significant for our understanding of (un)gram\-maticality feelings is that it suggests a profound disconnect between the evolution of language and the evolution of dependency sensitivity. Rather than evolving specifically for linguistic purposes, our feeling of ungrammaticality when encountering a malformed sentence may be an \textsc{exaptation}~-- the repurposing of a cognitive ability that evolved for entirely different reasons.

Like the feeling of disgust that evolved to help us avoid pathogens but now extends to moral violations, our sense of grammaticality likely built upon pre-existing pattern-recognition systems that long predated human language. Our hominid ancestors may have used these skills for parsing complex environmental patterns, recognizing social structures, or even processing the temporal sequences in conspecific vocalizations, long before the emergence of syntactic language.

This perspective helps explain why our feelings of (un)grammaticality can be so context-dependent and variable. The neural circuitry underlying these feelings wasn't optimized specifically for language but was co-opted from more general-purpose cognitive mechanisms that had evolved for other adaptive functions. The mechanism~-- detecting violations of expected patterns~-- served our ancestors well in many domains of life, and language simply hijacked this pre-existing system.

The squirrel monkey study also emphasizes how these mechanisms are inherently predictive. The monkeys weren't simply memorizing specific sequences; they were extracting abstract patterns and using them to make predictions about novel stimuli. When these predictions were violated, the monkeys showed increased attention~-- the animal equivalent, perhaps, of our feeling of ungrammaticality. This aligns with our discussion of prediction error as a core component of grammaticality intuitions.

If a squirrel monkey can detect a violation in a tone sequence following the same abstract pattern as Hungarian vowel harmony, then our own feelings about well-formedness in language likely arose from deep cognitive foundations shared across primates. Our human language faculty didn't create sensitivity to dependencies; it inherited and specialized mechanisms that were already millions of years old.

This evolutionary perspective challenges the notion that grammar is inherently linguistic. Rather, grammar~-- and our feelings about it~-- may be better understood as the application of domain-general pattern-recognition abilities to the specific realm of language. What we experience as a feeling of ungrammaticality may be nothing more than the same surprise signal our ancestors felt when encountering any unexpected pattern in their environment, now directed at linguistic stimuli.

\subsection{The social role of ungrammaticality feelings}

If our sense of ungrammaticality arose before human language, what purpose might it have served? Evidence from non-human primates and other social mammals offers compelling clues. The studies on squirrel monkeys by \citet{Ravignani2013} along with others on rock hyraxes by \citet{Kershenbaum2012} reveal that syntactic processing capabilities may have evolved for social purposes rather than for language per se.

Rock hyraxes look like scruffy rodents but are actually more closely related to elephants and manatees, a biological surprise package about the size of a sturdy rabbit. These charming, vocal mammals live in colonies among rocky outcrops across Africa and the Middle East, where they spend their days sunbathing, foraging, and engaging in complex social interactions punctuated by elaborate songs. That such an unassuming creature produces vocalizations with sophisticated syntactic structure challenges our assumptions about which animals might possess complex communication systems.

In squirrel monkeys, the ability to detect abstract dependencies between non-adjacent elements suggests a cognitive capacity for recognizing patterns in vocal sequences. Similarly, hyraxes demonstrate complex syntactic structures with geographical dialects that correlate with social proximity. Both findings point to a potential social function of syntax detection: group recognition and cohesion.

What we experience as a feeling of ungrammaticality might have originated as a social signal~-- an alert to potential out-group communication. For our primate ancestors, detecting ``incorrect'' vocal patterns could help identify unfamiliar individuals or groups, triggering appropriate caution or interest. This mechanism would be particularly valuable in environments where visual identification was limited, such as dense forests or during low-light conditions.

The correlation between vocal similarity and geographical distance in rock hyraxes suggests that syntactic patterns may serve as markers of social familiarity. As \citet{Kershenbaum2012} observed, song syntax similarity degrades with distance, but only within a limited range corresponding to typical dispersal patterns. This pattern is consistent with a social signaling system that evolved to distinguish local group members from strangers.

The mother-daughter vocal interactions in gibbons described by \citet{Koda2013} also suggest how syntactic patterns might be socially transmitted. Their finding that ``co-singing interactions between mothers and daughters may enhance vocal development'' suggests that social learning plays a crucial role in acquiring group-specific syntactic patterns.

Our human sense of ungrammaticality, then, may be an evolutionary repurposing of a much older social recognition system. What was once a mechanism for distinguishing in-group from out-group vocalizations has become a sophisticated tool for parsing linguistic structure. But the visceral nature of our reactions to ungrammatical sentences~-- that feeling of wrongness~-- may reflect this deep evolutionary history.

This evolutionary heritage may help explain the surprisingly emotional responses many people have to ``grammatical errors'' in modern contexts. Language policing~-- the practice of publicly correcting others' speech or writing~-- rarely stems from genuine communication needs. Rather, it often functions as a social gatekeeping mechanism. The grammar pedant who bristles at a split infinitive or corrects someone's \textit{whom} usage may be unconsciously engaging an ancient social marking system, one that evolved to identify who belongs and who doesn't.

The strong negative reactions to dialectal differences, casual speech patterns, or linguistic innovations reveal how deeply intertwined our sense of grammaticality is with social identity. When we feel that visceral discomfort at hearing someone \enquote{break the rules} of grammar, we may be experiencing the echo of an evolutionary adaptation that once helped our ancestors distinguish friend from stranger. In modern contexts, however, these instinctive reactions often serve to reinforce social hierarchies and exclude those who speak differently~-- suggesting that our feelings about grammaticality remain as much about social boundaries as they are about language itself.

\section{One intuition or many}
% TODO: source-grounding -- needs proper \citet{Knobe-forthcoming} and bib entry, or cut
Joshua Knobe (forthcoming from Ergo) argues that not only do different people have different intuitions but that individuals have conflicting intuitions.
